\section{Preliminaries: Halal Trust Context and Ethical Design Requirements}
\label{sec:prelim}

This section provides contextual background on halal certification trust gaps and derives \emph{engineering design requirements} from Islamic ethical principles.
It does not constitute a religious ruling (fatwa); rather, it informs system design analogously to value-sensitive design~\cite{friedman2017value}.

\subsection{Halal Certification Context}
Indonesia's halal certification ecosystem involves BPJPH registration and MUI auditing for many product categories~\cite{tieman2019halal}.
Reported issues include certificate forgery and inability of consumers to independently verify batch-level claims at point of purchase~\cite{bendhaou2021blockchain}.

\subsection{From Principles to Design Requirements}
Table~\ref{tab:ethics} maps selected principles to verifiable protocol properties.

\begin{table}[t]
  \caption{Mapping of ethical design requirements to HalalChain implementation.}
  \label{tab:ethics}
  \centering
  \begin{tabular}{@{}llp{0.35\linewidth}@{}}
    \toprule
    \textbf{Principle} & \textbf{Requirement} & \textbf{Implementation} \\
    \midrule
    Amanah & Immutable audit record & On-chain status transitions \\
    Thoyyib & Verified-only halal display & Status = Verified $\rightarrow$ consumer UI \\
    Adl & Distributed data access & Public chain + IPFS \\
    Hisba & Auditor accountability & \texttt{auditor} address + timestamp \\
    \bottomrule
  \end{tabular}
\end{table}

The current prototype grants \texttt{AUDITOR\_ROLE} via an administrator (OpenZeppelin \texttt{DEFAULT\_ADMIN\_ROLE}), which partially limits the \emph{Adl} property; Section~\ref{sec:discussion} discusses this limitation.
